\documentclass[notitlepage,onecolumn,showpacs,aps,prb]{revtex4-2}
\usepackage{geometry}
\usepackage[bottom]{footmisc}
\usepackage{color}
%\usepackage{doi}
\usepackage{hyperref}
%\hypersetup{hidelinks=true,}
%\usepackage[T1]{fontenc}
%\usepackage[utf8]{inputenc}
\usepackage{amssymb}
%\usepackage[nointegrals]{wasysym}
\usepackage{graphicx}
\usepackage{amsmath}
%\usepackage[breakable]{tcolorbox}
%\usepackage{esint}
\usepackage{epstopdf}
%\usepackage{color}
\usepackage{caption}
\usepackage{subcaption}
\usepackage{tensor}
\usepackage{slashed}
\usepackage{cancel}
%\usepackage{sidecap}
%\sidecaptionvpos{figure}{t}
\allowdisplaybreaks


\usepackage{feynmp-auto}
\setlength{\unitlength}{1mm}
\DeclareGraphicsRule{*}{mps}{*}{}

\newcommand{\bsym}[1]{\boldsymbol{#1}}
\newcommand{\mbf}[1]{\mathbf{#1}}
\newcommand{\op}[1]{\hat{#1}}
\newcommand{\grad}{\boldsymbol{\nabla}}
\newcommand{\Tr}{\text{Tr}}
\newcommand{\tr}{\text{tr}}
\newcommand{\N}{\mathcal{N}}
\renewcommand{\k}{\mathbf{k}}
\newcommand{\q}{\mathbf{q}}
\renewcommand{\r}{\mathbf{r}}
\renewcommand{\d}{\mathfrak{D}}
\newcommand{\F}{\mathcal{F}}
\renewcommand{\L}{\mathcal{L}}
\newcommand{\beq}{\begin{equation}\begin{aligned}}
\newcommand{\eeq}{\end{aligned}\end{equation}}

%\newcommand{\beq}{\begin{align}}
%\newcommand{\eeq}{\end{align}}

%\def\beq#1\eeq{\begin{align}#1\end{align}}

\newcommand{\todo}[1]{{\color{red} {\bf [ToDo: #1]}}}
\newcommand{\mv}[1]{{\color{green}{#1}}}

\DeclareMathOperator*{\sumint}{%
\mathchoice%
  {\ooalign{$\displaystyle\sum$\cr\hidewidth$\displaystyle\int$\hidewidth\cr}}
  {\ooalign{\raisebox{.14\height}{\scalebox{.7}{$\textstyle\sum$}}\cr\hidewidth$\textstyle\int$\hidewidth\cr}}
  {\ooalign{\raisebox{.2\height}{\scalebox{.6}{$\scriptstyle\sum$}}\cr$\scriptstyle\int$\cr}}
  {\ooalign{\raisebox{.2\height}{\scalebox{.6}{$\scriptstyle\sum$}}\cr$\scriptstyle\int$\cr}}
}

\graphicspath{{./}{./figs/}}


%\begin{figure}
%  \includegraphics[width=0.8\columnwidth]{filename}
%   \caption{fig caption}
%   \label{labelname}
%\end{figure}
  
\begin{document}

\title{The generalization of the non-relativistic Luttinger-Zeeman action to four spatial dimensions}
\date{\today}
\begin{abstract}
We derive the non-relativistic Luttinger-Zeeman action in arbitrary dimensions.
\end{abstract}
\author{Santanu Dey}

\maketitle

\section{Introduction}
\label{intro}

Let's consider the $j=3/2$ spin-orbit coupled Luttinger Hamiltonian in the presence of
an external magnetic field in $d=3$. In the interest of a higher dimensional extension, we 
write down a Hamiltonian
\beq
H&=\frac{4}{2(d-1)}\left[\left(\vec p\cdot\vec S\right)^2-\frac{j(j+1)}{d}p^2\right]+g \vec B\cdot\vec S, \text{ with},\\
\left[S_i,S_j\right]&=i\epsilon_{ijl}S_k \text{ for } d = 3.
\eeq
The Hamiltonian acts on $j=3/2$ fermionic spinors.
SO(3) rotation of the three-dimensional Euclidean space determines the form of the Zeeman coupling in
the Hamiltonian. In a field-theoretic language, the momentum-space Lagrangian
corresponding to this non-relativistic
Hamiltonian can be written as
\beq
\mathcal L =\psi^\dagger\left[\partial_\tau+H\right]\psi,
\eeq
where $\psi$ is a four-component SU(2) spinor with $j=3/2$. Under an infinitesimal 
SO(3) rotation, we have the following
transformation:
\beq
\substack{R\in {\rm SO(3)},\\ U(R) \in {\rm SU(2)}}: &\begin{cases}
  B_i\mapsto R_{ij}B_j, \\
  \psi^\dagger S_i\psi \mapsto \psi^\dagger U^\dagger(R)S_i U(R)\psi =R_{ij}\psi^\dagger S_j\psi
\end{cases}
\eeq
This demonstrates that the Zeeman vertex $\vec B\cdot\vec S$ is rotationally invariant. That the 
$j=3/2$ spin matrices $\vec S$ also transform like vector under SU(2) transformation $U$ 
follows from the Lie-algebraic isomorphism between the two groups
\beq
\rm Lie[SU(2)] = Lie[SO(3)],
\eeq
and it can be proved at the level of infinitesimal transformation by
exploiting the algebraic structures of the two groups.

The extension of the Zeeman coupling to higher dimension $d$ hinges on the isomorphism between the 
Lie algebras of the rotation group
SO(d) and the corresponding complex spinorial group Spin(d). For $d=3$, we have, 
Lie[Spin(3)] = Lie[SU(2)] = Lie[SO(3)]. For $d=4$, however, SO(4) has 6 generators with a Lie algebra,
\beq
[J_i,J_j]&=i\epsilon_{ijk}J_k,\\ 
[J_i,K_k]&=i\epsilon_{ijk}K_k,\\
[K_i,K_j]&=i\epsilon_{ijk}J_k,
\eeq
which can be decoupled with the parametrization $J_i^1=(J_i+ K_i)/2$ and
$J_i^2=(J_i- K_i)/2$,
\beq
[J^1_i,J^1_j]&=i\epsilon_{ijk}J^1_k,\\
[J^2_i,J^2_j]&=i\epsilon_{ijk}J^2_k,\\
[J^1_i,J^2_j]&=0,
\eeq
into two SU(2) Lie algebras. In other words \footnote{As exponentiation of Lie algebras generate Lie
groups we can say $\rm SU(2)\otimes SU(2) \rightarrow SO(4)$. In fact, this is a double cover,
$\rm SO(4)\simeq SU(2)\otimes SU(2)/\mathbb Z_2$
},
\beq
\rm Lie[SO(4)] = Lie[Spin(4)]=Lie[SU(2)]\oplus Lie[SU(2)].
\eeq
The spinor representation of SO(4) is therefore characterized by the $j$ quantum numbers of the 
two SU(2) groups, labeled by their spins $(j_1,j_2)$ with a combined highest weight $j_1+j_2$. 
First, we need to define spinors with total spin 
$j=3/2$ in the four dimensional space. This can be accomplished by choosing between 
$(3/2,0)$, $(1,1/2)$, $(1/2,1)$, and $(0,3/2)$ representations. By choosing one of the
two half-integer, fully-fermionic representations,
\beq
(\frac{3}{2},0)&: \vec J_1=\vec S, \vec J_2 = 0,\\
(0,\frac{3}{2})&: \vec J_1=0, \vec J_2 = \vec S,
\eeq
we can define four-component $j=3/2$ spinor fields. But this is not enough to obtain a rotationally
invariant theory of fermionic spinors coupled together with bosonic vector fields. The spin one
bosonic fields $B_i$
transforms as SO(4) vectors with the six generators of rotation 
$[\mathcal J^{ij}]_{kl}=i(\delta_{ik}\delta_{jl}-\delta_{il}\delta_{jk})$,
\beq
B_k\mapsto\left(\delta_{kl} 
-\frac{i}{2}\theta_{ij}[\mathcal J^{ij}]_{kl}\right)B_l.
\eeq
with infinitesimal rotation generated by the antisymmetric matrix $\theta$. The six
rotation generators can be compactly combined as,
\beq
\mathcal J=\begin{pmatrix}
  0&&K_1&&K_2&&K_3\\
  -K_1&&0&&J_3&&-J_2\\
  -K_2&&-J_3&&0&&J_1\\
  -K_3&&J_2&&-J_1&&0
\end{pmatrix},
\eeq
with the generators $J_i$ and $K_i$ following the algebra mentioned above. In the $(j_1,j_2)$ 
spinor basis, they are respectively spin $j_1+j_2$ and spin $j_1-j_2$ matrices.
Evidently, in a 
$(j,0)$ spinor basis, distinct SO(4) 
rotations generated by, e.g., $\mathcal J_{12}$ and $\mathcal J_{34}$ 
leads to the same transformation on the spinor field, as the latter transforms 
under the superposed SU(2) generator $\vec J_1=(\vec J+\vec K)/2$. 
We shall, therefore, need Dirac spinors in the $(3/2,0)\oplus (0,3/2)$ representation 
that transform uniquely under rotation.

Now, we have a $2\times (2\times 3/2+1)=8$ dimensional spinor representation in a four dimensional
space, and we wish write down an invariant scalar involving the spinor and tensor fields. For
ordinary Dirac spinors in $(1/2,0)\oplus (0,1/2)$ representation, this is achieved through 
generators of the Euclidean Clifford algebra in the spinor space,
\beq
\left\{\Gamma^i,\Gamma^j\right\}&=2\delta^{ij}, \\ 
\Gamma^{ij}&=\frac{i}{2}\left[\Gamma^i,\Gamma^j\right], S^{ij}=\frac{1}{2}\Gamma^{ij}.
\eeq
The matrices $S^{ij}$ generates SO(d) rotation
\beq
[\Gamma^i,S^{kl}]&=\tensor{[\mathcal J^{kl}]}{^i_j}\Gamma^j,\\
\left[S^{ij},S^{kl}\right]
&=i\left(\delta^{jk}S^{il}-\delta^{ik}S^{jl}
-\delta^{jl}S^{ik}
+\delta^{il}S^{jk}
\right).
\eeq
Note, that in the Euclidean space, raising and lowering indices do not introduce any sign factors. Using the 
first commutation relation, it
is now straightforward to show that under unitary transformation $U=\exp\left(-(i/2)\theta_{ij}S^{ij}\right)$
the Clifford matrices transform as, 
\beq
U^\dagger\Gamma^i U = \tensor{\exp\left(-\frac{i}{2}\theta_{kl}\mathcal{J}^{kl}\right)}
{^i_j}\Gamma^j = \tensor{R}{^i_j}\Gamma^j.
\eeq
This implies that under $SO(d)$ rotation generated by $d(d-1)/2$ distinct $\Gamma^{ij}$ matrices, 
the $(1/2,1/2)$ spinor bilinear product $\psi^\dagger\Gamma^i\psi$ transform as vectors. Similarly,
the product $\psi^\dagger\Gamma^{ij}\psi$ transform as rank-2 tensors. And a Zeeman coupling can 
be expressed as the invariant product 
\beq
\psi^\dagger\Gamma^{ij}\psi F_{ij}.
\eeq
In $d=2$ and $3$, the product reduces to the familiar Zeeman coupling. 

Now, dimensionality of the spinor space also restrict the number of generators of the Clifford group. For
a 2-dimensional (Weyl) spinor in the representation $(1/2,0)$,  there are three $2\times 2$ generators 
given by the
Pauli matrices
\beq
\Gamma^1=\sigma^1, \Gamma^2=\sigma^2, \Gamma^3=-i\sigma^1\sigma^2=\sigma^3,
\eeq
with $\Gamma^{ij}$ generating SO(3) rotation.
For the 4-dimensional $(1/2,0)\oplus(0,1/2)$ Dirac representation, five $4\times 4$ Clifford matrices can
be constructed by doing a Kronecker product of the $2\times 2$ Clifford matrices with $\sigma^1$ and
adding a component $\sigma^2\otimes 1$ \cite{ramond10},
\beq
&\Gamma^1=\sigma^1\otimes\sigma^1, \Gamma^2=\sigma^1\otimes\sigma^2,\\
&\Gamma^3=\sigma^1\otimes\sigma^3, \Gamma^4=\sigma^2\otimes 1,\\
&\Gamma^5=(-i)^2\Gamma^1\Gamma^2\Gamma^3\Gamma^4=\sigma^3\otimes 1,
\eeq
that constitutes an SO(5) Clifford algebra. The last Clifford matrix is obtain by,
$\Gamma^{2n+1}=
(-i)^n\Gamma^1\Gamma^2\cdots\Gamma^{2n}$. The construction reveals that 
$2n+1$ Clifford matrices are needed to span a spinor space of $2^n$ dimension.

For our $2^3=8$ dimensional $(3/2,0)\oplus (0,3/2)$ Dirac spinors, we shall need a construction with 
$7$ Clifford matrices. The construction may used to generate SO(7) rotation or two distinct 
SO(6) rotations with a (chiral) projection operator $P_{\pm}=(1\pm\Gamma^7)/2$ acting 
on the spinor space, splitting the 8 dimensional Dirac 
representation into two irreducible 4 dimensional Weyl representation $(3/2,0)$ and 
$(0,3/2)$. By sticking to the Dirac representation,
a SO(7) symmetric scalar can be written down in the eight-dimensional space $Z^a$,
\beq
\psi^\dagger\Gamma^{ab}\psi F_{ab}.
\eeq
Now, a reduction in dimensionality can be performed by declaring that $F_{ab}(Z)=F_{ab}(x)$,
where $x^i$ is the four-dimensional space of our interest. By the virtue of 
$F_{ab}(x)=\partial_{Z_a}A_b(x)-\partial_{Z_b}A_a(x)$ being the field strength tensor, 
the Zeeman coupling is uniquely determined.

\section{SO(9) Clifford algebra and how Murakami et al. \cite{murakami04} writes the Zeeman term}

Following Murakami et al. \cite{murakami04} and 
Janssen et al.'s \cite{janssen15} parametrization, the Clifford matrices are
defined as,
\beq
\Gamma^a= \sqrt\frac{1}{2d(d-1)}\Lambda^a_{ij}\{S^i,S^j\} \text{ with }
\Lambda^a_{ij}\Lambda^a_{kl}=\delta_{ik}\delta_{jl}+\delta_{il}\delta_{jk}-\frac{2}{d}\delta_{ij}\delta_{kl},
\eeq
where the second formula is due to Janssen et al. who consider a $d$ dimensional generalization
of the problem from $d=3$. $S_i$ are the spin $j/2$ angular-momentum matrices satisfying the
SU(2) algebra. Now, let us introduce the functions, $d_a(\vec p)=\sqrt\frac{d}{2(d-1)}\Lambda_a^{ij}p_i p_j$. 
From the above property of the $\Lambda^a$ matrices, it 
follows that,
\beq
\Gamma^a d_a(\vec p)
&=\frac{1}{2(d-1)}\Lambda^a_{ij}\Lambda_a^{kl}\{S_i,S_j\}p_k p_l,\\
&=\frac{4}{2(d-1)}
\left[\left(\vec p\cdot\vec S\right)^2-\frac{j(j+1)}{d}p^2\right],\\
&=H
\eeq
We are interested in enumerating the rotation generators $\Gamma^{ab}=\frac{i}{2}[\Gamma^a,\Gamma^b]$ 
and their relation to the $j=3/2$ spin operator $S^i=\eta^i_{ab}\Gamma^{ab}$. We invert the relation,
\beq
&\tr[S^i\Gamma^{cd}]=\eta^i_{ab}\tr[\Gamma^{ab},\Gamma^{cd}],\\
&\tr[S^i\Gamma^{cd}]=-d_\Gamma\eta^i_{ab}\left(\delta^{ad}\delta^{bc}-\delta^{ac}\delta^{bd}\right),\\
&\eta^{cd}_i=\frac{1}{2d_\Gamma}\tr[S_i,\Gamma^{cd}],
\eeq
where $d_\Gamma$ is the matrix dimension of the Clifford matrices and we have used the antisymmetry
properties of the $\eta^i_{ab}$ symbol to arrive at the result. Now, by substituting the definition of
$\Gamma^a$ matrices in terms of the spin matrices we have,
\beq
\eta_i^{ab}=\frac{i}{2d_\Gamma}\frac{1}{2d(d-1)}
\Lambda^a_{mn}\Lambda^b_{rs}\tr\left[S^i[\{S^m,S^n\},\{S^r,S^s\}]\right].
\eeq
The Zeeman coupling can now be written down as,
\beq
&\psi^\dagger S^i B_i\psi,\\
=&\psi^\dagger\Gamma^{ab}\psi\frac{\eta^{ab}_i\epsilon_{ijk}}{2}F^{jk},\\
=&
g^{ab}_{jk}
\psi^\dagger\Gamma^{ab}\psi F^{jk}.
\eeq
In the current computation, it is this $g^{ab}_{jk}$ symbol that I computed using the $d$ dimensional
formulae for $\Lambda^a_{ij}$ and $\delta_{ij}$.

\section{Janssen et al.'s \cite{janssen15} way of doing higher dimensional extension}

With $(d+2)(d-1)/2$ number of $\Gamma^a$ matrices, 
Janssen et al. writes down a general quadratic band touching (QBT) Hamiltonian, 
\beq
H=G_{ij}p_i p_j \text{ with } G_{ij}=\sqrt\frac{d}{2(d-1)}\Lambda^a_{ij}\Gamma_a,
\eeq
in general $d$ dimensions. From the earlier discussion, we can surmise that in $d=4$, 
the nine $16\times 16$ Clifford matrices can be combined to construct a sixteen-dimensional 
SO(9) representation. Now, let us think about the Fermioic spinor space for a moment. We have
moved from three-dimensional Euclidean space with $SO(3)$ rotational invariance to
four-dimensional Euclidean space with $SO(4)$ rotational symmetry. Furthermore, we wish
to have a quantized description of the fermionic fields in their spinor representation. 
In the three-dimensional space, the label $j$ of the SU(2) spinors describe the spin
of the homomorphic representation SO(3). With $\rm SO(4)=SU(2)\otimes SU(2)/\mathbb Z_2$, 
the notion of spin is associated with the angular momentum of its SO(3) subalgebra. The
decomposition goes as follows (see Sec.~\ref{intro} for explanation):

\begin{table}[htb!]
\begin{tabular}{l*{6}{c}r}
Dimensions             & 2 & 4 & 4 & 8 & 16  & 16\\
\hline
  Rep. of $\rm SU(2)\otimes SU(2)$ 
  & (1/2,0) 
  & (1/2,1/2) 
  & $\rm (1/2,0)\oplus (0,1/2)$ 
  & $\rm (3/2,0)\oplus (0,3/2)$ 
  & (15/2,0)
  & $\rm (7/2,0)\oplus (0,7/2)$ 
  \\
  Rep. of $\rm SO(3)$ sub-algebra & 1/2 
  & $\rm 1\oplus 0$ 
  & 1/2 
  & 3/2 
  & 15/2
  & 7/2
  \\
\end{tabular}
\end{table}
Clearly, a half-integer spin description of the QBT Hamiltonian involves fermions of 
$j>3/2$ angular momentum, with the identification of the angular momentum operators 
as the generator of the SO(3) subalgebra of the SO(4) spatial symmetry. 
However, this characterization 
holds onto the canonical notion of angular momentum and spin tied to a SO(3) submanifold 
of SO(4). In principle, we could dispense with this construction and notions of subspace spins
altogether. We could instead identify combinations of 
the six generators $\vec J^1$ and $\vec J^2$ 
of a $(j_1,j_2)$ representation of $\rm SU(2)\otimes SU(2)$ as the six generators 
$\vec J=(\vec J_1+\vec J_2)/2$, $\vec K=(\vec J_1-\vec J_2)/2$ of the SO(4) rotation. In this 
way, the sixteen-dimensional fermionic spinor space can be described with a
$(j_1,j_2)=(3/2,3/2)$ Weyl representation (see Sec.~\ref{intro} for the notation used). This is
inherently a multi-integer spin representation of the SO(4) rotation group,
\beq
(3/2,3/2)=3\oplus 2\oplus 1\oplus 0\oplus 0,
\eeq
where the last zero comes from the $\vec K$ representation of SO(4)/SO(3). However, these
are artificial spinors aiding us towards a dimensional increment of the three-dimensional
$j=3/2$ QBT fermions. We wish to write down a Zeeman like coupling in four-dimensional
space that reduces to a regular three-dimensional Zeeman term after dimensional reduction.
The obvious candidate is the term,
\beq
H_Z = \psi^\dagger \Sigma_{ij} \psi F_{ij},
\eeq
where $F_{ij}$ and $\Sigma_{ij}$ are second rank tensors under SO(4) rotation to form
this invariant coupling. Under the rotation of the SO(4) vector space, we must have,
\beq
\substack{R\in {\rm SO(4)},\\ U(R) \in {\rm SU(2)\otimes SU(2)}} : &\begin{cases}
  F_{ij}\mapsto R_{ik}R_{jl}F_{kl}, \\
  \psi^\dagger \Sigma_{ij}\psi \mapsto \psi^\dagger U^\dagger(R)\Sigma_{ij} U(R)\psi 
  =R_{ik}R_{jl}\psi^\dagger \Sigma_{kl}\psi
\end{cases}
\eeq
SO(4) spatial rotations in 
Janssen et al.'s construction with a large sixteen-dimensional spinor space leaves the
quadratic band touching tensor $G_{ij}$ invariant. For an infinitesimal rotation, 
$R_{ij}=\delta_{ij}-(i/2)\theta_{kl}[\mathcal J^{kl}]_{ij}$ and 
$ U_{\alpha\beta}=\delta_{\alpha\beta}-(i/2)\theta_{kl}[S^{kl}]_{\alpha\beta}$,
the construction yields,
\beq
U^\dagger(R)G_{ij}U(R) &= R_{ik}R_{jl}G_{kl}\\
G_{ij}-\frac{i}{2}\theta_{uv}[G_{ij},S^{uv}]
&=
G_{ij}-\frac{i}{2}\theta_{uv}
\left(
[\mathcal J^{uv}]_{ik}\delta_{jl}
+\delta_{ik}[\mathcal J^{uv}]_{jl}
\right)
G_{kl}.
\eeq
The sixteen-dimensional spinor space is spanned by the matrices 
$1$, $\Gamma^a$, and $\Gamma^{ab}$, so we can decompose $S^{uv}=A^{uv}+B^{uv}_a\Gamma^a
+C^{uv}_{ab}\Gamma^{ab}
$. In order to derive these $A$, $B$, and $C$ symbols, we compute the commutator in the left
hand side and find that,
\beq
[G_{ij},S^{uv}]=-(2\sqrt 2i)\sqrt\frac{d}{d-1}
\left[C^{uv}_{ab}\Lambda^b_{ij}\Gamma^a+B^{uv}_a\Lambda^{ij}_e\Gamma^{ea}\right].
\eeq
It is clear that for $S^{uv}$ to generate SO(4) rotation, the symbols $B^{uv}_a$ must identically 
vanish. For the right hand side, we have the expression,
\beq
&\left(
[\mathcal J^{uv}]_{ik}\delta_{jl}
+\delta_{ik}[\mathcal J^{uv}]_{jl}
\right)
G_{kl}\\
=&-i\sqrt\frac{d}{2(d-1)}\Gamma^e\left(\delta_{iv}\Lambda^e_{ju}
-\delta_{iu}\Lambda^e_{vj}+\Lambda^e_{iu}\delta_{jv}-\Lambda^e_{iv}\delta_{ju}\right).
\eeq
Comparing both sides we have the relation for the $C$ symbols,
\beq
4C^{uv}_{ab}\Lambda^b_{ij}
=\left(\delta_{iv}\Lambda^a_{ju}
-\delta_{iu}\Lambda^a_{vj}+\Lambda^a_{iu}\delta_{jv}-\Lambda^a_{iv}\delta_{ju}\right),
\eeq
from which the antisymmetry of the symbol under $u\leftrightarrow v$ is evident, as expected. To
enumerate the symbol, we contract both sides with the $\Lambda^c_{ji}$, and use the property 
$\tr\left[\Lambda^a\Lambda^b\right]=2\delta^{ab}$ \cite{janssen15} to obtain,
\beq
C^{uv}_{ab}=\frac{1}{4}\left(
\Lambda^a_{ui}\Lambda^b_{iv}
-
\Lambda^a_{vi}\Lambda^b_{iu}
\right)
=\frac{1}{4}[\Lambda^a,\Lambda^b]_{uv},
\eeq
which completes our mapping of the SO(4) rotation generator in this higher-dimensional
spinor space
\beq
S^{ij}&=C^{ij}_{ab}\Gamma^{ab},\\
&=\frac{i}{4}\left(
\left[\Gamma^a\Lambda^{ik}_a\right]
\left[\Gamma^a\Lambda^{kj}_a\right]
-
\left[\Gamma^a\Lambda^{jk}
_a\right]\left[\Gamma^a\Lambda^{ki}_a\right]
\right)
\eeq
In an analogy to the lower-dimensional Zeeman coupling (the spatial part of the Pauli term
of a spatio-temporal electrodynamics) we now write down an expression for a higher-dimensional
Zeeman Hamiltonian,
\beq
H_Z=\frac{1}{2}\psi^\dagger\Gamma^{ab}\psi C^{ij}_{ab}F_{ij}
\eeq


\bibliography{lqed}

\end{document}
